% Section 1 -- Introduction
% Draft prose, written 9 November 2026 from 01-introduction-outline.tex
% (frame "A+B"; the outline stays on disk and is no longer typeset). One
% paragraph per outline point, in the outline's order. Every number is taken
% from the text or tables of Sections 2, 4, 5 and 6: Section 5 numbers use
% 65--75, Section 6 numbers 65--74. \open{} marks a decision still to take,
% \addref{} a reference not yet in epc-pd.bib.
\section{Introduction}\label{sec:intro}

% \input{sections-pd6/01-introduction-outline}

% 1. The prediction, and a fact that seems to confirm it
One of the basic predictions of the theory of investment is that people invest less
when fewer years remain to collect the return
\citep{becker1962investment,ben1967production}. For instance, a longer expected life raises investment in schooling
\citep{jayachandran2009life,oster2013limited,fortson2011mortality}, and
workers who expect to retire later train more
\citep{montizaan2010effects}.\footnote{The prediction also holds outside
human capital: executives in their final years in office cut R\&D
\citep{dechow1991executive}. There is a debate about which horizon matters:
Americans worked fewer lifetime hours as they lived longer, so the
remaining lifetime cannot be the whole story
\citep{hazan2009longevity,cervellati2013life}.} Schooling and training stay
with the worker, who can neither buy them ready-made nor sell them on. However, when we look at durable goods like housing, which change hands over
time, investments are embodied in the good and pass from seller to buyer.
The energy efficiency of a home is one such investment. Efficiency
improvements lower the energy bill for years and stay with the house when
its owner sells, so a buyer can acquire them with the home. At the
same time,  mortality risk rises over the lifecycle,  which by the
standard argument shortens the horizon (although older owners, who move
less, stay longer in their homes). In this paper we study  which of these forces -- a horizon that shortens with age
and efficiency that is bought and sold with the house -- shapes
the energy efficiency of  homes over the lifecycle of households. In the rest of the paper, for brevity, we call the first force the horizon channel and the second the purchase window, or simply the window.


% 2. The English facts
We first use survey and register data in England to document three stylized facts. First, older owners live in less efficient homes. In the English Housing
Survey, among owner-occupied homes built by 1990, 21 percent of those
owned by households aged 25--34 are in Energy Performance Certificate (EPC, henceforth) bands
A--C, against 9.7 percent of those owned by households aged 65--75, representing a gap of about 11 percentage points.  Second,  following a home across successive certificates in the
EPC register, we show that  efficiency is not fixed: each year about 7
percent of homes built by 1990 in bands D--G move to A--C, and about 15 percent
of those in A--C move down. Third, efficiency is tied to when the
owner bought. Owners who bought recently live in about equally
efficient homes at every age.




% 3. Two routes to efficiency, and which force acts on which
 What do these three stylized facts say about the two forces discussed above? The answer depends on
how an owner gets an efficient home. An owner can make her home efficient in two ways: (i) buy a home that is
already efficient, or (ii) renovate the one she has. The difference between the two lies in what a
later sale returns. In the first case, a buyer pays for efficiency in the price and gets it
back when she sells, if the home is still efficient. In the second case,   a renovator pays for the work, and a sale  returns only
part of it.\footnote{Prices capitalise the value of energy savings
\citep{myers2019home,aydin2020capitalization,fuerst2015energy}, and
retrofits typically cost more than the savings they deliver
\citep{fowlie2018weatherization}.}  The horizon channel acts on  the owner's renovation decision. An older owner has fewer years to recover her investment, so
she might renovate less. On the other hand,  the purchase window affects whether an
owner can get efficiency through a purchase. Because a buyer gets back at resale most of what she paid for
efficiency, buying an efficient home is the cheaper route. However, it is open
only when an owner moves, and older owners in Britain move less often
\citep{banks2012housing}. Between moves, efficiency wears off, as the
second fact shows. An owner who stays therefore lives in a home that
becomes less efficient each year, and older owners have usually stayed the
longest since their last purchase. So, both the horizon channel and the purchase window
can produce the first fact.  The third fact favours the window: among owners who bought
recently, older and younger owners live in about equally efficient homes,
so the gap opens with the time since purchase rather than with age itself.
It does not settle the question, because under the horizon channel, too,
an owner who stops renovating sees her home's efficiency fall with the
time since purchase.




% 4. Why a model, and the model (old paragraph 4 parked, entry 10)
Why does it matter whether the declining age gradient comes from the
horizon channel or the purchase window? The two channels have different
implications for an ageing population. If the horizon channel determined
the gradient, then an increase in life expectancy would raise investment in
efficiency. However, if the purchase window determined it, longer
lives would do little to raise investment: what matters is how often owners buy, and
older owners buy less often. This means, from a policy perspective, knowledge of the underlying mechanism is crucial to determine whether to raise efficiency by incentivizing renovation or by acting at the point of sale.



To separate the implication of the horizon channel and the purchase window, we need two quantities that are not fully observable: (i) how fast efficiency wears off relative to how fast mortality
rises, and (ii) how much cheaper efficiency is at a purchase than through a
renovation.\footnote{The register gives the rate at which homes lose a band,
but only for homes that are certified again, and a change in band is not
always a physical change. The register matched to Land Registry prices gives
the price premium and the proportional cost of a renovation, but not the fixed
cost of the work. The model we build below estimates both the rate at which
efficiency is lost and the fixed cost (Section~\ref{sec:identification}).}  To obtain them, we build a
lifecycle model of households with
income risk, mortality risk, and borrowing limits, built on the standard
ingredients of lifecycle models of housing
\citep{cocco2005portfolio,yang2009consumption,kaplan2020housing}. Each year
a household chooses consumption and liquid saving, whether to rent or own,
the size of her home, and her mortgage, and, if she owns, whether her home
is efficient or not. An owner keeps her home, renovates
it in place, or sells; a seller or a renter rents, buys an efficient home,
buys an inefficient one, or buys an inefficient one and renovates it at
once. Efficiency lowers the energy bill, gives
comfort, and sells at a price premium, and each year it is lost with a
fixed probability. Renovating has a fixed and a variable cost, buying and
selling have transaction costs, and a bequest motive values what heirs
receive, including the premium on an efficient home. The model holds
tastes, information, and the cost of the work fixed across ages.\footnote{Other explanations could also produce the facts: older owners
may value efficiency less, know less about it, find the work more
burdensome, or prefer to run down their home equity. These are preference-based explanations and we do not rule them out. However, we ask whether the dynamics of the investment alone can
account for the gradient.} Any age gradient in the model must therefore come from the horizon channel or from the purchase window, which opens when an owner buys and so depends on transaction costs, borrowing limits, bequests, and income.




% 5. What the model says before any numbers
The model has many discrete and continuous choices
and occasionally binding borrowing constraints, so we
proceed in two steps. In the first, we simplify the model to
isolate the key mechanisms analytically. This gives three results. (i) An
owner is less willing to renovate the home she lives in the more likely
she is to leave it, whether by moving or by dying, because a sale
returns only part of the cost of the work. This is the horizon channel.
(ii) A buyer's willingness to choose an efficient home depends
little on how long she expects to stay, because she, or her heirs, recover
most of the premium when the home is sold, provided it is still
efficient. (iii) Renovating is therefore worthwhile for fewer owners than
buying an efficient home is, because it costs more once what a sale returns
is counted. Further, this difference grows with the chance of leaving. Some
owners would not renovate the home they live in, yet would choose an
efficient home if they were buying; whether they live in one depends on
whether they move. This is the purchase window. These results give the
direction of each channel. How much each contributes to the age gradient is
the second step, for which we solve and estimate the full model.


% 6. Solving the model
Solving and estimating this model is cumbersome. An owner with a
mortgage has five continuous states --- liquid wealth, the mortgage balance,
two components of income, and the size of the home --- two home types, and up
to 25 remaining mortgage terms, and she chooses among up to six alternatives
with two or three continuous choices in each. This means a grid-based solution method cannot be employed.\footnote{A grid with 50 points per
continuous state would have up to $1.6\times10^{10}$ points at each age.} We instead approximate the value and policy
functions at each age with neural networks, trained backward from the
oldest age on randomly drawn states
\citep{maliar2021deep,azinovic2022deep,druedahl2026deep}. Two devices make
this work. Taste shocks over the alternatives smooth the kinks that
discrete choices put in the value function, so that a network can fit it
\citep{iskhakov2017endogenous}. And the parameters we estimate are inputs of
the networks, so one training covers the whole estimation search
\citep{druedahl2026nonconvex}.


% 7. Estimation and the test the model passes untargeted
We estimate seven parameters by the simulated method of moments, targeting
eleven moments on ownership, wealth, consumption, and the share of efficient homes by
age. The estimates deliver the two quantities that the facts alone could
not. First, renovating a home costs between 2.7 and 4.2 times the extra
outlay of buying the same home efficient, depending on its
size.\footnote{The fixed cost of renovating is identified by how steeply
the efficient share falls with age, because older owners rely more on
renovating in place (Section~\ref{sec:identification}).} Second, efficiency
wears off fast: an efficient home loses its efficiency with a probability
of 16 percent a year, close to the 15 percent of efficient homes that drop
to bands D--G each year in the register, a rate we do not
target.\footnote{The targeted moments pin down only a combination of how
long efficiency lasts and how much it is worth each year, because both
raise the efficient share at every age; the register's rate separates
them. Part of that rate reflects reassessment rather than physical change
\citep{crawley2019quantifying}, so we also report the main results when
efficiency lasts 25 years instead of six.} The estimated model reproduces
10 percentage points of the 11 percentage point fall in the efficient share between ages
25--34 and 65--75 (Figure~\ref{fig:moment-fit}).



% 8. Results 1 and 2: the window, not the horizon (old paragraphs 8-9 parked, entry 11)
Using the estimated model we first untangle whether the age gradient comes from the
purchase window or from the horizon channel. The horizon channel turns out to be weak because the horizon itself hardly differs between young and old
owners. What matters for a renovation is not how long the owner will live,
but how long she will benefit from the work. If efficiency lasted forever,
an owner aged 25 could expect to benefit for 27 discounted years and an owner aged 75
for 10. But efficiency wears off long before most owners die, so the
benefit lasts about five years at 25 and about four at 75. However, the purchase window
differs sharply with age. Owners of every age make similar choices when
they act. A little over two in five buyers choose an efficient home, and
fewer than two in a hundred owners of inefficient homes renovate in a
year. What falls with age is how often owners buy, and this fall accounts
for almost all of the decline with age in the number of owners who gain an
efficient home. Older owners do not
renovate instead, because of the cost of the work, which a buyer of
an efficient home avoids. When we remove the option to buy an efficient
home, the gap between the youngest and the oldest owners falls from about
ten percentage points to about one.\footnote{The result survives a much
slower loss of efficiency. If efficiency lasts 25 years rather than six,
the horizon now plays a role, but purchases still account for about 90
percent of the decline.}

% 10. Result 3: which policies reach older owners
The United Kingdom aims to bring as many
homes as possible to band C by 2035 \citep{hmg2017cleangrowth}.
Households headed by someone aged 65 or over own more than a third of
owner-occupied homes in England \citep{mhclg2024ehs}, and their number is
projected to grow by about a fifth in the ten years to 2032, twice as fast
as households overall \citep{ons2025households}. We use the estimated model for a series of counterfactuals, focusing on policies which governments generally apply to improve the efficiency of homes. Because the gradient
comes from the window, the route through which a policy acts decides whom
it reaches. We compare steady states in partial equilibrium and report each
policy as a change from the model without it. A
subsidy that halves the variable cost of renovating raises the efficient
share by about 8 points among owners aged 65--74 but by less than 3 points
at 25--34, so it narrows the gap between them by more than 5 points, at a
cost of about \pounds290 per household a year. It reaches older owners because they
are the ones who rely on renovating in place. A point-of-sale standard,
which requires every home that is bought to be efficient, raises the share
the most, by about 13 points among all owners, but by more among the young
(about 16 points) than the old (about 10), so it widens the gap by about 6
points. It also slows turnover. A household that would have bought an
inefficient home now pays more to buy, and some put off moving: owners move
and buy about a quarter less often. Each move that does not happen is a
window that does not open, so the standard gives back part of its own
gain. Lowering the costs of buying and selling works in the opposite
direction. When we halve them, owners buy more than twice as often, more
windows open, and the efficient share rises by about 5 points among all
owners. But owners who expect to move sooner renovate less: the rate of
renovation among owners who stay falls by about a third.

% 10b. Result 3, continued: changing the payoff instead of the cost
The counterfactual policies we discussed above change what it costs to get an efficient home. A policy
can instead change what an efficient home pays its owner: a lower energy
bill each year, and a higher price when the home is sold. We vary each in
turn. Doubling energy prices, as a tax on household energy would, raises
the efficient share in about the same proportion among young and old
owners, by about 40 percent. It lifts the whole age profile without
changing its shape, but it lowers homeownership by about 3 percentage points. The price premium works differently, and against a
natural conjecture. If older owners renovate little because the market pays
little for efficiency, a larger premium should bring them into efficient
homes. In the model, doubling the premium does raise renovation, most among
older owners, because a sale then returns more of the cost of the work. But
a buyer pays the premium, so buying an efficient home becomes more
expensive and fewer buyers choose one. Young owners, who get most of their
efficiency by buying, lose the most. The age gap almost closes, but because
the homes of young owners become less efficient, not because those of older
owners improve.



% 11. Why it matters beyond homes: quality is chosen at the transaction
The mechanism is not specific to energy. Much of the quality of a durable
is chosen when it is bought, because that is when changing it is cheap.
Models of durable adjustment explain when households replace what they own
\citep{grossman1990asset,eberly1994adjustment,attanasio2000consumer}. Our
results add that the quality a household holds then depends on how often
it transacts, so a household that rarely transacts holds quality chosen
long ago. The fuel economy of the car fleet improves as old cars are
scrapped and new ones bought, so a regulation that raises the price of new
cars delays the improvement
\citep{gruenspecht1982differentiated,stavins2006vintage,jacobsen2015vehicle},
and when owners replace or resell a car is itself an economic choice
\citep{adda2000balladurette,stolyarov2002turnover}.
Homes differ from cars in two ways. First, a home is rarely replaced but often
resold. Second, its efficiency can also be raised in place. The moment of
gaining efficiency is still mostly the purchase, and how often homes change hands depends on
the age of their owners and on the costs of transacting. Older owners
therefore gain efficient homes more slowly than younger ones,
even if no owner's incentive to renovate changes with age. Building codes
reach only new homes \citep{levinson2016energy,jacobsen2013building}, and
most of the homes that will stand in 2035 are already built.\footnote{England
had 25.6 million dwellings in March 2024 and added about 210,000 in the
following year \citep{mhclg2025supply}. At that rate, homes already built
will make up more than nine-tenths of the stock in 2035.}

% 12. Contribution relative to the literature
We contribute to four literatures. The first studies how the horizon shapes
investment \citep{becker1962investment,ben1967production} and finds that a
longer life raises schooling
\citep{jayachandran2009life,oster2013limited,fortson2011mortality}, though
which horizon counts is still debated
\citep{hazan2009longevity,cervellati2013life}. We answer this question for
an investment embodied in a durable that wears off and is left behind at a
move: its payoff horizon is short at every age, so the remaining lifetime is
not what sets the age profile of investment; the timing of transactions is. The second asks
why households appear to underinvest in energy efficiency
\citep{jaffe1994energy,allcott2012energy,gerarden2017assessing,fowlie2018weatherization},
and often points to high implicit discount rates or heterogeneity in them
\citep{hausman1979individual,newell2015individual,bento2012there}. Low
efficiency among older owners need not reflect high discount rates or
misperceived savings; in our model it reflects when the cheap moment to
invest arrives. Buyers pay a premium for a better certificate
\citep{fuerst2015energy}, and ratings bunch just above band thresholds
\citep{collins2018bunching}, so the certified band carries value of its
own; we take the band as the object of choice and its premium as part of
the payoff. The third builds lifecycle models of housing with tenure, home
size, mortgages, and moving
\citep{cocco2005portfolio,yang2009consumption,attanasio2012modelling,kaplan2020housing}
and studies saving and housing in old age
\citep{denardi2010elderly,davidoff2010home,nakajima2017reverse,nakajima2020home}.
We add a quality attribute of
the home that can be bought with it or added in place, and that wears off.
The fourth shows that transaction costs reduce mobility
\citep{vanommeren2005new,hilber2017transfer,best2018housing}.\footnote{A mortgage locked in at a low rate also reduces mobility \citep{quigley1987interest,ferreira2010housing,fonseca2024mortgage}. The model has no such lock-in, but it keeps the mortgage rules under which a borrower who moves is more restricted than one without a mortgage (Section~\ref{sec:model-environment}). Like transaction costs, this makes moving less attractive for borrowers, and a household that moves less often has fewer chances to buy an efficient home. We would therefore expect the same channel to apply.} We
show that transaction costs, through mobility, also slow
the upgrading of the quality of the housing stock. Three papers are closest
to ours. \citet{fleischer2026energy} solve the renovation problem of a Dutch
household from age 60 that stays in its home, and find that mortality risk
and liquidity constraints delay renovation; our owners can also buy an
efficient home, and it is this second route, not the shorter horizon, that
separates young owners from old. The two findings on the horizon do not
conflict. In their model a renovation lowers energy use for good, so its
benefit lasts as long as the household lives and mortality cuts it short. In
ours efficiency wears off in about six years, so the benefit is short at
every age and the remaining lifetime changes it little.
\citet{droes2024home} show in Dutch survey
data that households with higher incomes both improve their homes more often
and sort into homes that are already efficient; we study the same two
margins by age, in a model that explains why older owners rely less on
buying and how policy acts on each margin. \citet{reusens2026renovation}
estimate that the Flemish obligation on buyers of the least efficient homes
to renovate within five years lowered the price of these homes by about 2
percent; we simulate a stricter standard, under which every home that is
bought must be efficient, and trace which owners it reaches and how it
changes how often they move. Methodologically, we solve a large
lifecycle model with neural networks, building on \citet{maliar2021deep},
\citet{azinovic2022deep}, and
\citet{druedahl2026deep,druedahl2026nonconvex}.



% 14. Road map
Section~\ref{sec:data} describes the data and three facts.
Section~\ref{sec:model} sets out the model and the three propositions.
Section~\ref{sec:estimation} describes the solution and the estimation.
Section~\ref{sec:results} separates the horizon from the window, and
Section~\ref{sec:policy} studies policies. Section~\ref{sec:conclusion}
concludes.
